\documentclass[a4paper,11pt]{article}
\usepackage[margin=0.6in]{geometry}
\usepackage{enumitem}
\usepackage{titlesec}
\usepackage{hyperref}
\usepackage{fontawesome}
\usepackage{array}
\usepackage{xcolor}
\usepackage{setspace}
\usepackage{nopageno}
\definecolor{linkcolor}{RGB}{0, 0, 255}
\titleformat{\section}{\large\bfseries\centering}{\thesection}{0em}{}[\titlerule]
\titlespacing{\section}{0pt}{5pt}{3pt}
\setlength{\parindent}{0pt}
\setstretch{0.82}
\setlist[itemize]{leftmargin=0.4cm, itemsep=0pt, parsep=0pt, topsep=1pt}
\hypersetup{colorlinks=true, linkcolor=linkcolor, urlcolor=linkcolor}
\newcommand{\resumeItem}[1]{\item[\circ] #1}
\newcommand{\linktext}[2]{{\color{linkcolor}\href{#1}{#2}}}
\begin{document}
\begin{minipage}[t]{0.5\textwidth}
    \textbf{\Large SURAJ PRASAD KALAUNI}\\
    \href{https://www.linkedin.com/in/surajprasaddtu}{LinkedIn} \\
    \href{https://github.com/surajprasad32}{GitHub}
\end{minipage}
\begin{minipage}[t]{0.5\textwidth}
    \begin{flushright}
        Email: \href{mailto:spkalauni7789@gmail.com}{spkalauni7789@gmail.com}\\
        Mobile: +91 8267081462
    \end{flushright}
\end{minipage}
\vspace{-0.1cm}
\section*{EDUCATION}
\begin{tabular}{@{}p{0.7\textwidth}p{0.3\textwidth}@{}}
    Delhi Technological University & Delhi, India \\
    \textit{M.Tech Software Engineering, GPA: 8.33} & \textit{Aug 2024 - Aug 2026} \\
    Dr. A.P.J Abdul Kalam Technological University & Lucknow, India \\
    \textit{B.Tech Computer Science and Engineering, GPA: 8.12} & \textit{Jun 2020 - Mar 2024} \\
\end{tabular}
\vspace{-0.1cm}
\section*{SKILLS SUMMARY}
\begin{itemize}
\item \textbf{Languages}: Python, C++, C, Java, Golang
\vspace{0.05cm}
\item \textbf{ML Frameworks \& Inference Tools}: PyTorch, Hugging Face
Transformers, TensorFlow, vLLM, BitsAndBytes, ONNX, LangChain, SGLang
\vspace{0.05cm}
\item \textbf{Inference \& Optimization}: Speculative Decoding, INT8/FP16
Quantization, Structured Pruning \& Compression, Sparse Attention,
KV-Cache Optimization, Post-Training Techniques, RAG, Prompt Engineering
\vspace{0.05cm}
\item \textbf{Systems \& Infrastructure}: Docker, Kubernetes, AWS (EC2, S3,
Lambda), Linux, REST APIs, CI/CD, Observability \& Monitoring
\item \textbf{Core Computer Science}: Data Structures \& Algorithms, Operating
Systems, Computer Networks, System Design, OOP, DBMS
\end{itemize}
\vspace{-0.1cm}
\section*{WORK EXPERIENCE}
\textbf{Samsung Research and Development $|$ Software Engineer Intern}
\hfill July 2025 -- January 2026
\begin{itemize}
    \item Engineered LLM-powered agentic workflows integrated into internal
    developer tools; applied prompt chaining and tool-use patterns to automate
    multi-step developer tasks, cutting manual overhead significantly.
    \item Optimised transformer-based inference pipelines for production
    deployment — applied batching, prompt compression, and caching strategies,
    reducing average inference latency by \textbf{15\%} on internal benchmarks.
    \item Built and containerised full-stack services (Next.js, TypeScript)
    using Docker \& Kubernetes on AWS, achieving \textbf{99.9\%} uptime with
    improved system observability and monitoring.
    \item Collaborated with cross-functional ML and platform teams to validate
    model outputs, debug inference failures, and tune generation parameters
    across multi-agent orchestration frameworks.
\end{itemize}
\vspace{-0.1cm}
\section*{PROJECTS}
\textbf{Speculative Decoding — LLM Inference Acceleration}\hfill 2025
\begin{itemize}
    \item Implemented speculative decoding from scratch using GPT-2 as draft
    model and GPT-2 XL as target verifier; reproduced token acceptance and
    rejection mechanism as described in Leviathan et al. (2023).
    \item Measured \textbf{2.1$\times$ end-to-end throughput speedup} over
    standard autoregressive decoding at equivalent output quality; analysed
    acceptance rate vs. speedup tradeoff across temperature settings.
    \item Profiled token generation latency, KV-cache memory usage, and batch
    throughput across sequence lengths to quantify performance gains on CPU
    and single-GPU environments.
\end{itemize}
\textbf{LLM Inference Benchmarking \& Quantization Pipeline}\hfill 2025
\begin{itemize}
    \item Built end-to-end benchmarking tool for open-source LLMs
    (LLaMA 3.2-1B, Mistral-7B) measuring tokens/second, per-token latency,
    peak memory, and batch throughput under varying load.
    \item Applied INT8 post-training quantization via BitsAndBytes; achieved
    \textbf{38\% reduction in GPU memory} with less than 1.2\% perplexity
    degradation on WikiText-2 compared to FP16 baseline.
    \item Exported quantized models to ONNX format and validated inference
    outputs, documenting latency-accuracy tradeoffs across precision levels
    (FP32, FP16, INT8) for reproducible evaluation.
\end{itemize}
\textbf{Transformer Pruning \& Compression for Efficient Inference}\hfill 2024
\begin{itemize}
    \item Designed structured pruning pipeline for DistilBERT; applied
    magnitude-based attention-head pruning at 10\%, 20\%, and 30\% sparsity
    levels and measured downstream accuracy on SST-2 classification.
    \item Achieved \textbf{27\% reduction in inference latency} at 20\%
    sparsity with only 0.8\% accuracy drop; benchmarked against knowledge
    distillation baseline to compare compression strategies.
    \item Analysed sparsity-accuracy tradeoff curves and visualised attention
    head importance scores to identify redundant heads — findings directly
    informed pruning threshold selection.
\end{itemize}
\vspace{-0.1cm}
\section*{PUBLICATIONS}
\begin{itemize}
    \item Retinal Fundus Image and Deep Learning: A Novel Approach for
    Cardiovascular Risk Prediction, Published in AIP Conference Proceedings,
    2025.
    \item Introduced a novel method for cardiovascular risk prediction,
    utilizing retinal fundus images and deep learning techniques, achieving
    79.41\% accuracy in classifying cardiovascular risk.
    \item DOI: \linktext{https://doi.org/10.1063/5.0247009}{https://doi.org/10.1063/5.0247009}
\end{itemize}
\vspace{-0.1cm}
\section*{ACHIEVEMENTS}
\begin{itemize}
    \item \linktext{https://drive.google.com/file/d/119HcGj2dlvhP-H1cSROeO51gGnfPOQP5/view}{GATE CSE 2024}: Qualified with Marks: 40.48, Score: 490/1000, Rank: 5943 among 123967.
    \item \linktext{https://drive.google.com/file/d/1G_rEs7Rse5fOASrs1HE2DaDU_To_P00e/view}{GATE DA 2024}: Qualified with Marks: 38, Score: 362/1000, Rank: 5973 among 39210.
\end{itemize}
\end{document}
